\documentclass[10pt,a4paper]{amsart}
\usepackage[margin=19mm]{geometry}
\usepackage{amsmath,amssymb,mathtools}
\usepackage[scr=rsfs,cal=euler]{mathalfa}
\usepackage{xfrac}
\usepackage[unicode,hidelinks,bookmarks=false]{hyperref}
\newcommand{\ie}{i.e.\ }
\newcommand{\cf}{cf.\ }
\newcommand{\resp}{respectively\ }
\newcommand{\Subsection}{\subsection}
\providecommand{\nobreakdash}{\nobreak\hbox{-}}
\setlength{\emergencystretch}{2em}
\multlinegap0pt
\usepackage{amssymb}
\usepackage{xcolor}
\usepackage{tikz}
\newtheorem{proposition}{Proposition}
\title{Pansu pullback and spectral complexes}
\author{Filippa Lo Biundo, Francesca Tripaldi}
\date{Source: arXiv:2603.22134v2 (CC BY 4.0)}
\begin{document}
\maketitle

\noindent\emph{Source and licence.} Pansu pullback and spectral complexes. Version arXiv:2603.22134v2 (CC BY 4.0), distributed by arXiv under the Creative Commons Attribution 4.0 International licence, http://creativecommons.org/licenses/by/4.0/, as recorded in the arXiv licence metadata for this exact version and shown on the abstract page https://arxiv.org/abs/2603.22134v2. Full licence notice: \texttt{licenses/arxiv-2603.22134-CC-BY-4.0.txt}.\par\smallskip

\noindent\textit{Editorial scope.} Counted unit: the article's proposition prop:triangularity-classical-differential (source0.tex lines 1025--1086). The article's complete original proof of it is delivered with it (source0.tex lines 1088--1190, one \texttt{proof} environment of 103 lines). No mathematics is written, repaired or re-derived: the statement and the proof are the article's own lines, in the article's own notation, and no retained line is substituted or re-flowed.  No further source-local context is needed: every cross-reference of every retained line resolves inside the two delivered spans.  Nothing was omitted from any retained span, so the package carries no omission record.  No retained line contains a citation, so the wrapper carries no bibliography and no bibliography entry.  No retained line contains a figure environment, an image-inclusion command or drawing code, so no image is delivered and the package declares no asset dependency.  Editorial formatting only, in addition: an AMS \texttt{amsart} preamble that restates the article's own declarations and macros used by the retained lines, namely the environments \texttt{proposition} on separate counters, together with the standard packages those lines need.  The retained environments print in this document as Proposition 1 in order of appearance, whereas the article prints the same items with the numbering of its own full document.  The complete native source of the article as distributed in the arXiv e-print for this exact version is bundled unchanged as \texttt{sources/196999/source0.tex}.
\begin{proposition}[The classical differential of a Pansu differentiable map]
\label{prop:triangularity-classical-differential}
Let \(G_1\) and \(G_2\) be Carnot groups with stratified Lie algebras
\[
\mathfrak g_1=V_1^1\oplus\cdots\oplus V_{s_1}^1
\ \text{ and }\ 
\mathfrak g_2=V_1^2\oplus\cdots\oplus V_{s_2}^2\,.
\]
Let $\phi\colon\mathbb R^{n_1}\longrightarrow\mathbb R^{n_2}$
be the expression in exponential coordinates of a Euclidean
\(\mathcal C^1\) map
\(\varphi:G_1\to G_2\) and assume \(\varphi\) is Pansu
differentiable at $\exp(x)\in G_1$.

Set
\[
\widetilde\phi
:=
L_{-\phi(x)}\circ\phi\circ L_x
\]
and let
\[
a_{i,j}
:=
\Pi_i\bigl(d(\widetilde\phi)_0(e_j)\bigr),
\]
where \(e_j\) belongs to the $l_j^1$-th layer of $\mathfrak{g}_1$, while
the \(i\)-th coordinate in the target belongs to the $l_i^2$-th layer.

Then
\[
a_{i,j}=0
\qquad\text{whenever}\qquad
l_i^2>l_j^1\,.
\]
Moreover,
\[
\Pi_i\bigl(D_P\phi(x)(e_j)\bigr)
=
\begin{cases}
a_{i,j}\ , & l_i^2=l_j^1\,,\\[2mm]
0\ ,       & l_i^2\neq l_j^1\,.
\end{cases}
\]
Equivalently, the classical differential in left-translated
coordinates satisfies
\[
d(\widetilde\phi)_0(V_q^1)
\subseteq
\bigoplus_{p\leq q}V_p^2\ \text{ for every }q=1,\ldots,s_1\,,
\] and
\[
D_P\phi(x)\big|_{V_q^1}
=
\operatorname{pr}_{V_q^2}
\circ d(\widetilde\phi)_0\big|_{V_q^1}.
\]
Thus, with respect to the bases adapted to the stratifications,
the matrix representing \(d(\widetilde\phi)_0\) is block upper
triangular, and the Pansu derivative \(D_P\phi(x)\) is given
precisely by the block diagonal part.
\end{proposition}

\begin{proof}
Since $\widetilde\phi(0)=0$ and we are assuming \(\widetilde\phi\) is Euclidean differentiable at \(0\), for every
\(e_j\) we have
\[
\widetilde\phi(\lambda^{l_j^1}e_j)
=
\lambda^{l_j^1}d(\widetilde\phi)_0(e_j)
+
o(\lambda^{l_j^1})
\ \text{ as }\ \lambda\to0^+.
\]
Taking the \(i\)-th coordinate gives
\begin{equation}
\label{eq:euclidean-expansion-coordinate}
\Pi_i\bigl(\widetilde\phi(\lambda^{l_j^1}e_j)\bigr)
=
\lambda^{l_j^1}a_{i,j}
+
o(\lambda^{l_j^1})\,.
\end{equation}
On the other hand, by Pansu differentiability,
\[
\delta_{1/\lambda}\circ\widetilde\phi\circ\delta_\lambda(e_j)
\longrightarrow
D_P\phi(x)(e_j)\ 
\text{ as }\ \lambda\to0^+.
\]
Since \(e_j\) has weight \(l_j^1\), we have $\delta_\lambda(e_j)=\lambda^{l_j^1}e_j$.
Moreover, the dilation of \(G_2\) acts on the \(i\)-th coordinate by
multiplication by \(\lambda^{l_i^2}\). Consequently,
\begin{equation}
\label{eq:pansu-component-limit}
\lambda^{-l_i^2}
\Pi_i\bigl(\widetilde\phi(\lambda^{l_j^1}e_j)\bigr)
\longrightarrow
\Pi_i\bigl(D_P\phi(x)(e_j)\bigr)\,.
\end{equation}
The Pansu derivative commutes with the homogeneous dilations and
therefore preserves the stratifications, that is $D_P\phi(x)(V_q^1)\subseteq V_q^2$.
It follows that
\begin{equation}
\label{eq:off-diagonal-pansu}
\Pi_i\bigl(D_P\phi(x)(e_j)\bigr)=0
\ \text{ whenever }\ 
l_i^2\neq l_j^1\,.
\end{equation}
One should then distinguish the following three cases.

Suppose first that \(l_i^2>l_j^1\). By
\eqref{eq:pansu-component-limit} and
\eqref{eq:off-diagonal-pansu},
\[
\Pi_i\bigl(\widetilde\phi(\lambda^{l_j^1}e_j)\bigr)
=
o(\lambda^{l_i^2})\,.
\]
Combining this with
\eqref{eq:euclidean-expansion-coordinate}, we obtain $\lambda^{l_j^1}a_{i,j}+o(\lambda^{l_j^1})
=
o(\lambda^{l_i^2})$, so when
dividing by \(\lambda^{l_j^1}\) we get
\[
a_{i,j}+o(1)
=
o\bigl(\lambda^{l_i^2-l_j^1}\bigr)\,.
\]
Since \(l_i^2-l_j^1>0\), letting \(\lambda\to0^+\) yields $a_{i,j}=0$.

Suppose next that \(l_i^2=l_j^1\). Dividing
\eqref{eq:euclidean-expansion-coordinate} by \(\lambda^{l_i^2}\), we find $
\lambda^{-l_i^2}
\Pi_i\bigl(\widetilde\phi(\lambda^{l_j^1}e_j)\bigr)
=
a_{i,j}+o(1)$.
Comparing this with \eqref{eq:pansu-component-limit} gives
\[
\Pi_i\bigl(D_P\phi(x)(e_j)\bigr)=a_{i,j}\,.
\]

Finally, if \(l_i^2<l_j^1\), then
\eqref{eq:euclidean-expansion-coordinate} gives $\lambda^{-l_i^2}
\Pi_i\bigl(\widetilde\phi(\lambda^{l_j^1}e_j)\bigr)
=
\lambda^{l_j^1-l_i^2}\bigl(a_{i,j}+o(1)\bigr)$.
Since \(l_j^1-l_i^2>0\), the right-hand side converges to zero.
Therefore,
\[
\Pi_i\bigl(D_P\phi(x)(e_j)\bigr)=0\,.
\]
Notice that, unlike the case when $l_i^2>l_j^1$, no condition is imposed on \(a_{i,j}\).

We have therefore proved that $a_{i,j}=0$ whenever $l_i^2>l_j^1$, and that the Pansu derivative retains precisely the entries for which $l_i^2=l_j^1$. This is equivalent to
\[
d(\widetilde\phi)_0(V_q^1)
\subseteq
\bigoplus_{p\leq q}V_p^2
\ \text{ and }\ 
D_P\phi(x)\big|_{V_q^1}
=
\operatorname{pr}_{V_q^2}
\circ d(\widetilde\phi)_0\big|_{V_q^1}\,.
\]
\end{proof}

\end{document}
